\section{Conclusions}\label{sec:conclusions}

The published description of the numerical-relativity particle-in-cell method in Ref.~\cite{Wu:2026nrpic} suffices for an independent implementation once one unstated local-frame choice and open test settings are supplied. That implementation reproduces the single-particle geodesic benchmarks to all published digits and recovers every self-gravitating test in kind. Ten thousand macroparticles, evolved in minutes to hours on one graphics card, capture black hole formation, cluster stability, and the collapse, inspiral, and merger of two black holes inside a cluster.

A discrete trajectory pusher that updates particle momentum in a local orthonormal frame requires an explicit tetrad gauge. Midpoint evaluation of the Boris rotation across a finite time step leaves a second-order truncation error that depends on how that local frame rotates along the trajectory. A multi-pass grid filter on deposited matter sources requires a sharpening compensation pass in self-gravitating evolutions at low particle number. Plain binomial passes widen every particle cloud like gravitational softening and weaken the collective potential well enough to make stable clusters bounce outward. Cubic-lattice initial data at low particle number must rescale macroparticle masses in the outermost surface layer to match the continuum rest mass. The count of lattice sites inside a sphere differs from its continuum volume in cells and otherwise shifts the mass of the black hole formed in collapse.

Apparent-horizon mass measurements require an independent check from a geometric mass diagnostic evaluated on the spatial slice alone. On a resolved vacuum black hole, comparison with the slice Misner--Sharp mass exposed unfilled derivative ghost zones across mesh-block faces in the host code's apparent-horizon finder. Cloud-in-cell deposition should place macroparticle clouds on the finest local mesh level wherever adaptive mesh refinement remains bounded. Level-capped coarser clouds trade higher constraint residuals and shifted apparent-horizon masses near a forming black hole for bounded deep refinement on random samples. In a relativistic cluster with two collapsing overdensities, initial data must specify a mutually consistent host-envelope radius and orbital speed to fix the post-merger outcome. Through envelope gravity and accretion, the speed-matched radius and larger stated radius bracket the published evolution between complete post-merger envelope capture and a wide black hole orbit.

Among the effects that remain undetermined for future work, what initiates the coherent radial core oscillation that grows and saturates in unperturbed stable clusters at this particle number? Why do inward-kicked unstable clusters collapse later and the two black holes inside a cluster merge earlier than in Ref.~\cite{Wu:2026nrpic}? Would initial data that solve the constraints for superposed cluster and overdensity matter remove the post-merger apparent-horizon mass excess above the total mass? How does the apparent-horizon mass deficit after prolate collapse divide between macroparticle count and cell spacing across the axial caustic stage?

Language-model agents under human direction built the coupled solver, independent reference solutions, and diagnostics from the published method text alone. Predictions written before dedicated control runs separated true causes from guesses across tests in which several initial explanations failed. A method description that an independent builder can follow and a benchmark size that fits on one graphics card make such predictive checks cheap to repeat.
